\section{Experiment}
\label{sec:experiment}

\subsection{Setup}

We present key experiment configurations here, with detailed setups in Appendix~\ref{sec:appendix/exp_setup}.

\xhdr{Baselines}
We compare diverse methods under equal parameter budgets, using Qwen3-\{0.6B,1.7B,4B\}-Base~\citep{yang2025qwen3} as backbones.
We compare \name over the following baselines:
(1) \emph{Standard}, which always verbalizes directly and reduces to the standard Qwen model;
(2) \emph{SoftThink}, a latent optimization method, implemented following the official design~\citep{zhang2025soft} on top of the Standard model;
(3) \emph{AlwaysThink}, similar to \name, but always iterates twice for all tokens during training and inference;
(4) \emph{Ouro}, a looped transformer that can scale iteration depths, implemented following the official design~\citep{zhu2025ouro}.
Unless otherwise specified, \name, Ouro, and AlwaysThink all use a maximum of two iterations, fine-tuned from the same Qwen3 backbone on the same training data and recipe.
Ouro reaches its highest performance when set to maximum iterations, so we report results under this setup.
% We also compare with dynamic query routing via matrix factorization~\citep{ong2024routellm}, routing between MobileLLM-R1-360M~\citep{zhao2025mobilellm} and Qwen3-1.7B, as well as between Qwen3-0.6B and Qwen3-4B, to match average parameter sizes of 0.6B and 1.7B.
% \todo{mention Ouro is the best performing version, since it drops accuracy if iterating soon}
% For DiT~\citep{kim2025learning}, we follow the official training script, replacing the training data with the Open-R1 dataset.
% \todo{Add query-level routing setup}

\xhdr{\name setup}
To match the total parameter count of \name (with LoRA and decider) with that of baselines, we prune one layer from the LLM backbone before training.
The layer is chosen to minimize the increase in validation loss. 
We also report results for an unpruned variant, TaH+, which adds less than 3\% extra parameters from LoRA and decider. 
The detailed parameter composition is shown in Table~\ref{tab:param_composition}. 
Following~\citep{fu2025r2r}, we set the continuation threshold $c_{\text{threshold}}=0.9$ with about 7\% of tokens being iterated twice.
The oracle policy $\pi$ uses Qwen3-0.6B, 1.7B, and 4B as reference models, respectively.

\xhdr{Training scheme}
All models are trained on the balanced Open-R1~\citep{openr1} mixture (math, QA, and code; 100K samples) using supervised fine-tuning. 
To fit memory and compute limits, we exclude responses longer than 8,192 tokens; 4B models additionally truncate at 4,096 tokens; all other training settings follow the official Open-R1 script.
% The filtered dataset contains 480M tokens, with 1\% reserved for validation. 
Each method is sufficiently trained for 5 epochs, and we select the checkpoint with the lowest validation loss as the final model. 
All backbones are initialized from the corresponding Qwen3-Base model.


\phantomsection
\label{sec:eval_setup}
\xhdr{Evaluation setup}
We evaluate across challenging reasoning benchmarks, including GSM8K~\citep{cobbe2021training}, MATH500~\citep{hendrycks2021measuring}, AMC23 (American
Mathematics Competitions), AIME25 (American Invitational Mathematics Examination), OlympiadBench (denoted as Olympiad)~\citep{he2024olympiadbench}, MBPP++~\citep{mbpp}, HumanEval++ (denoted as HE++)~\citep{humaneval}, GPQA-Diamond (denoted as GPQA)~\citep{rein2023gpqa}, and MMLU-STEM (denoted as MMLU)~\citep{mmlu}.
The maximum generation length is set to 8,192 tokens for all benchmarks, except GSM8K and MMLU-STEM, which use 4,096 due to their simpler problems and larger size.
Performance is reported as pass@1 under a zero-shot CoT setting, using a sampling temperature of 0.6.
We generate one sample per problem for large datasets (MATH500, OlympiadBench, etc.), and eight samples per problem for small datasets (AMC23, AIME25).

\subsection{Performance}
\label{sec:experiment/performance}

\xhdr{Benchmark evaluation}
We validate \name's reasoning ability across all nine benchmarks.
Table~\ref{tab:performance_by_category_transposed} presents performance results for models at 0.6B, 1.7B, and 4B parameter sizes.
Compared with the strong Standard Qwen3 baselines, we observe that existing approaches (AlwaysThink,
SoftThink, and Ouro) yield only marginal improvements when fine-tuned from base.
In contrast, \name achieves consistent gains across benchmarks.
For 0.6B and 1.7B models, \name delivers average improvements of 3.0\% and 3.8\% over Standard, respectively; TaH+, which adds less than 3\% additional parameters, further pushes these gains to 5.3\% and 6.2\%.
Compared to the concurrent work Ouro, \name and TaH+ achieve 3.8-4.4\% and 6.1-6.8\% gains, respectively.
For 4B models, which are trained with a 4K context length due to resource limits, \name and TaH+ also achieve gains of 1.7\% and 2.2\%.

\begin{table*}[t]
    \caption{Accuracy comparison across benchmarks.
    Best results are highlighted in bold. $^*$4B models are trained with $\leq$4K lengths due to resource constraints; AlwaysThink is excluded due to Out-Of-Memory (OOM) during training.}
    \label{tab:performance_by_category_transposed}
    \centering
    \setlength\tabcolsep{4pt} % default value: 6pt
    \begin{tabular}{lcccccccccc}
    \toprule
    ~ &
    AIME25 & Olympiad & AMC23 & MATH500 & GSM8K & GPQA & MMLU & HE++ & MBPP++ & Average \\
    \midrule
    \multicolumn{11}{c}{\textit{0.6B}} \\
    \midrule
    Standard &
    1.9 & 15.4 & 22.7 & 39.9 & 58.2 & 31.1 & 54.2 & 16.8 & 28.8 & 29.9 \\
    SoftThink &
    \underline{2.9} & 14.0 & 22.2 & 39.6 & 55.9 & 24.7 & 53.0 & 14.3 & 29.5 & 28.5 \\
    Ouro &
    2.1 & 14.2 & 19.7 & 37.4 & 56.6 & \textbf{35.4} & 54.0 & 18.9 & 23.5 & 29.1 \\
    AlwaysThink &
    1.3 & 12.6 & 21.9 & 37.8 & 52.6 & 30.8 & 51.4 & 9.1 & 13.8 & 25.7 \\
    \rowcolor{tahgray}
    TaH &
    2.1 & \underline{19.1} & \underline{24.1} & \underline{46.2} &
    \underline{63.6} & 29.0 & \underline{56.4} & \underline{21.6} & \underline{33.9} & \underline{32.9}{\scriptsize$/+3.0$} \\
    \rowcolor{tahgray}
    TaH+ &
    \textbf{4.6} & \textbf{20.6} & \textbf{24.7} & \textbf{51.8} &
    \textbf{67.6} & \underline{31.3} & \textbf{59.0} & \textbf{22.0} & \textbf{35.1} & \textbf{35.2}{\scriptsize$/+5.3$} \\
    \midrule
    \multicolumn{11}{c}{\textit{1.7B}} \\
    \midrule
    Standard &
    10.8 & 33.8 & 39.7 & 67.8 & 80.2 & 30.3 & 74.1 & 39.0 & 51.9 & 47.5 \\
    SoftThink &
    5.4 & 30.7 & 40.3 & 64.8 & 80.0 & \underline{33.3} & 73.5 & 43.3 & 49.1 & 46.7 \\
    Ouro &
    10.8 & 31.3 & 40.6 & 68.2 & 79.8 & 32.8 & 72.4 & 40.9 & 45.5 & 46.9 \\
    AlwaysThink &
    7.5 & 31.0 & 40.9 & 63.2 & 74.2 & 30.5 & 69.6 & 16.4 & 25.6 & 39.9 \\
    \rowcolor{tahgray}
    TaH &
    \underline{13.8} & \underline{37.2} & \underline{40.9} &
    \underline{71.4} & \textbf{84.8} & \underline{33.3} &
    \underline{74.8} & \underline{50.0} & \underline{55.3} & \underline{51.3}{\scriptsize$/+3.8$} \\
    \rowcolor{tahgray}
    TaH+ &
    \textbf{15.4} & \textbf{37.6} & \textbf{48.4} &
    \textbf{72.6} & \underline{84.5} & \textbf{39.4} &
    \textbf{76.6} & \textbf{51.5} & \textbf{57.5} & \textbf{53.7}{\scriptsize$/+6.2$} \\
    \midrule
    \multicolumn{11}{c}{\textit{4B$^*$}} \\
    \midrule
    Standard &
    24.2 & 50.4 & 64.2 & 84.2 & \underline{91.7} &
    46.2 & 85.4 & 69.2 & 65.9 & 64.6 \\
    % AlwaysThink &
    % -- & -- & -- & -- & -- & -- & -- & -- & -- & -- \\
    SoftThink &
    25.8 & 50.2 & 63.1 & 85.0 & \textbf{92.5} & \textbf{50.3} & 86.0 & 69.2 & 66.7 & 65.4 \\
    Ouro &
    25.0 & \underline{51.4} & 64.4 & 83.8 & 90.7 & \underline{50.0} & 85.9 & \underline{70.7} & 66.7 & 65.4 \\
    \rowcolor{tahgray}
    TaH &
    \underline{27.1} & 50.5 & \textbf{69.7} & \textbf{85.8} & 91.0 & 48.5
     & \underline{86.2} & 70.1
     & \underline{67.7} & \underline{66.3}{\scriptsize$/+1.7$}\\
    \rowcolor{tahgray}
    TaH+ &
    \textbf{27.9} & \textbf{52.6} & \underline{68.1} & \underline{85.6} & \underline{91.7} &
    49.0 & \textbf{86.6} & \textbf{72.0} & \textbf{68.1} & \textbf{66.8}{\scriptsize$/+2.2$} \\
    \bottomrule
    \end{tabular}
\end{table*}


\xhdr{Hardware-agnostic efficiency}
Tables~\ref{tab:token_iter_stats} and~\ref{tab:efficiency_analysis_theo} report the average iteration depth, per-token FLOPs, and memory access of \name.
On average, \name performs 1.07 iterations per token. It significantly undercuts the 2.08-2.18$\times$ FLOPs and memory access of AlwaysThink, matching the overhead of Standard with only 4-5\% additional overhead. See Appendix~\ref{sec:appendix/exp/theo_eff} for more details.

\begin{table}[tb]
    \caption{Real-world decoding performance on a single A800 GPU, tested on AIME25 with 8K max token length.}
    \label{tab:real_world_perf}
    \centering
    \setlength\tabcolsep{4pt}
    \begin{tabular}{lcc>{\columncolor{tahgray}}c}
    \toprule
     & Standard & AlwaysThink & \name \\
    \midrule
    Avg. Depth & 1.00 & 2.00 & 1.06 \\
    Memory (GB)        & 4.3   & 6.8   & 4.6   \\
    Latency (s)        & 210.6 & 747.2 & 301.4 \\
    Throughput (token/s) & 38.9 & 11.0  & 27.2  \\
    \bottomrule
    \end{tabular}
\end{table}

\xhdr{Real-world efficiency}
We compare efficiency of 1.7B models at 8K length on a single NVIDIA A800 GPU.
As shown in Table~\ref{tab:real_world_perf}, \name iterates twice on only 6\% of tokens on AIME25, with 1.48$\times$ lower memory overhead and 2.48$\times$ faster decoding than AlwaysThink, while achieving higher accuracy.
More detailed efficiency experiment setup and runtime breakdown are shown in Appendix~\ref{sec:appendix/exp/real_world_efficiency}.
    
% \todo{add temperature=0 testing}


%%% Training dynamics %%%
% Figure
% X: training step
% Y: training loss and evaluation loss

\begin{figure}[tb]
    \centering
    \includegraphics[width=0.8\columnwidth]{figure/eval_loss_0_6.pdf}
    \caption{Training dynamics of the LLM backbone on Qwen3-0.6B-Base. \name converges rapidly and achieves lower perplexity.}
    \label{fig:training_llm}
\end{figure}

\begin{figure}[tb]
    \centering
    \includegraphics[width=0.8\columnwidth]{figure/iter_decider_acc.pdf}
    \caption{Iteration-decider accuracy vs. epoch (Qwen3-0.6B).}
    \label{fig:training_decider}
\end{figure}

\begin{figure}[tb]
    \centering
    \includegraphics[width=0.8\columnwidth]{figure/continuation_threshold.pdf}
    \caption{GSM8K accuracy with respect to continuation threshold. Numbers in brackets show the percentage of tokens iterated twice.}
    \label{fig:continuation_threshold}
\end{figure}


\xhdr{Training dynamics}
During stage~1 (LLM backbone training), \name performs iterations according to the oracle policy. As shown in Figure~\ref{fig:training_llm}, it converges notably faster than the Standard baseline and also achieves much lower validation perplexity.
During stage~2 (iteration-decider training), the neural decider successfully imitates the oracle strategy. It reaches about 83\% accuracy at predicting the oracle's iteration decisions, as shown in Figure~\ref{fig:training_decider}.


\subsection{Design Choice Exploration}
\label{sec:exp/dse}

We demonstrate the effectiveness and robustness of each design choice of \name through ablation studies.
All experiments in this section train \name and its variants on the math subset of Open-R1 and evaluate on MATH500, AMC23, and OlympiadBench.

%%% Architecture %%%
% \begin{table}[tb]
%     \caption{Ablation study on architecture designs, using 0.6B models. 
%     % \textit{Causal}: standard causal attention where tokens only attend to depth-1 contexts. 
%     % \textit{Duo-causal}: our proposed attention allowing tokens to attend to current and all shallower iterations.
%     }
%     % \vspace{-8pt}
%     \label{tab:ablation}
%     \centering
%     \begin{tabular}{lll|cccc}
%     \toprule
%     \textbf{LoRA}& \textbf{Residual}& \textbf{Attention} & \textbf{MATH500} & \textbf{AMC23}  & \textbf{Olympiadbench} & \textbf{Average}\\
%     \midrule
%     \cmark & \cmark & Duo-causal & 51.2 & \textbf{32.5}  & \textbf{23.9} & \textbf{35.9}\\
%     \midrule
%     \xmark & \cmark & \multirow{2}{*}{Duo-causal}  & \textbf{51.6} & 29.7  & 22.4& 34.6 {\scriptsize$/-$1.3}\\
%     \xmark & \xmark & ~ & 49.2 & 22.5  & 21.2& 31.0 {\scriptsize$/-$4.9}\\
%     \midrule
%     \cmark & \cmark & \textit{Causal} & 47.8 & 24.4  & 19.4 & 30.5 {\scriptsize$/-$5.4}\\
%     \bottomrule
%     \end{tabular}
% \end{table}

\begin{table}[tb]
\caption{Ablation study on design choices of TaH-0.6B. Each row varies one aspect from the \name configuration (marked with \colorbox{tahgray}{gray}).}
\label{tab:ablation}
\centering
\footnotesize
\setlength\tabcolsep{2pt}
\begin{tabular}{l@{\hskip 0pt}cccc}
    \toprule
    \textbf{Variant} & \textbf{MATH500} & \textbf{AMC23} & \textbf{Olympiad} & \textbf{Average} \\
    \midrule
    \rowcolor{tahgray}
    \quad \name & \textbf{51.2} & \textbf{32.5} & \textbf{23.9} & \textbf{35.9} \\
    \midrule
    \multicolumn{5}{c}{\textit{Model Architecture}} \\
    \midrule
    \multicolumn{5}{l}{\textit{Iteration Depth} (\name: \colorbox{tahgray}{Neural decider})} \\
    \quad Always-1 & 47.2 & 23.4 & 18.8 & 29.8{\scriptsize$/-$6.1} \\
    \quad Always-2 & 32.8 & 15.6 & 10.2 & 19.5{\scriptsize$/-$16.4} \\
    \addlinespace
    \multicolumn{5}{l}{\textit{Attention} (\name: \colorbox{tahgray}{Duo-causal})} \\
    \quad Causal@iter1 & 47.8 & 24.4 & 19.4 & 30.5{\scriptsize$/-$5.4} \\
    \quad Causal@current & 42.0 & 23.8 & 16.4 & 27.4{\scriptsize$/-$8.5} \\
    \addlinespace
    \multicolumn{5}{l}{\textit{Depth Adapters} (\name: \colorbox{tahgray}{LoRA + Residual})} \\
    \quad w/o LoRA & 51.6 & 29.7 & 22.4 & 34.6{\scriptsize$/-$1.3} \\
    \quad w/o LoRA \& Res. & 49.2 & 22.5 & 21.2 & 31.0{\scriptsize$/-$4.9} \\
    \midrule
    \multicolumn{5}{c}{\textit{Training Scheme}} \\
    \midrule
    \multicolumn{5}{l}{\textit{Supervision Type} (\name: \colorbox{tahgray}{Token-only})} \\
    \quad Token+latent & 49.4 & 29.6 & 15.9 & 31.6{\scriptsize$/-$4.3} \\
    \addlinespace
    \multicolumn{5}{l}{\textit{Iteration Policy during LLM training} (\name: \colorbox{tahgray}{Oracle policy})} \\
    \quad Decider-based & 44.8 & 24.1 & 17.3 & 28.7{\scriptsize$/-$7.2} \\
    \quad Dynamic & 11.0 & 2.8 & 2.7 & 5.5{\scriptsize$/-$30.4} \\
    \addlinespace
    \multicolumn{5}{l}{\textit{Discrepancy metric in $\pi$} (\name: \colorbox{tahgray}{Top1 mismatch})} \\
    \quad Cross-entropy & 47.4 & 21.2 & 20.4 & 29.7{\scriptsize$/-$6.2} \\
    \quad Entropy & 42.0 & 21.9 & 16.9 & 26.9{\scriptsize$/-$9.0} \\
    \bottomrule
\end{tabular}
\end{table}
    



\xhdr{Model architecture}
% We examine the effectiveness of two key architectural components: LoRA and residual connections for smooth objective transition, and the duo-causal attention pattern for leveraging extended context. 
(1) \textbf{Iteration Scheme}.
As shown in Table~\ref{tab:ablation}, decider-based iteration outperforms the \textit{Always-1} and \textit{Always-2} alternatives, confirming the practical benefits of selective iteration, even with imperfect decisions from the neural decider. Note that for \textit{Always-1}, duo-causal attention degenerates to regular causal attention.
(2) \textbf{Duo-Causal Attention}.
Replacing duo-causal attention with standard causal attention variants causes significant drops:
(a) attending only to the first iteration (Causal@iter1) drops by 5.4\%;
(b) attending only to the current iteration (Causal@current) drops even more, by 8.5\%.
They confirm duo-causal attention's essential role in cross-depth information flow.
(3) \textbf{Depth Adapters}.
Removing LoRA and residual connections leads to consistent drops, confirming their beneficial roles in objective transition across iterations.

% Figure
% Training scheme (iter decision)
\xhdr{Training scheme}
% We examine how supervision types and iteration behavior during backbone training affect performance. 
(1) \textbf{Supervision type.} 
Some previous work supervises all iteration depths with next-token labels, while \name only supervises final output tokens.
As shown in Table~\ref{tab:ablation}, such \textit{token+latent} supervision underperforms \name. This aligns with our intuition that different iterations should focus on their respective objectives.
(2) \textbf{Iteration policy during LLM training.}
We compare our static oracle strategy $\pi$ with two alternatives.
The \textit{decider-based} approach trains the iteration decider first, then uses it during backbone training. It suffers from the coupling challenge discussed in Section~\ref{sec:training}.
The \textit{dynamic} approach recalculates the oracle using the evolving backbone in Equation~\ref{eq:ours_oracle_depth}, facing the same coupling challenge and causing training collapse. 
These results support our selection of the static oracle policy.
(3) \textbf{Discrepancy metric in $\pi$.}
We compare three discrepancy metrics to trigger latent iteration in $\pi$: top-1 mismatch, entropy, and cross-entropy (detailed definitions in Appendix~\ref{sec:appendix/exp/oracle_label}).
As Table~\ref{tab:ablation} shows, top-1 mismatch yields the best result, confirming its empirical effectiveness.

\xhdr{Iteration label robustness}
We further evaluate how sensitive \name is to the quality of its oracle iteration labels.
For this evaluation, all \name-1.7B variants are trained on the math split while varying only how the iteration depth labels are constructed.
As shown in Table~\ref{tab:oracle_label_robustness}, labels from a smaller 0.6B reference and labels with 10\% random flips both still improve over Standard-1.7B.
The clean same-scale 1.7B reference gives the strongest average gain, so we use this label source in the main experiments.

\begin{table*}[t]
\caption{Robustness of \name-1.7B to imperfect oracle iteration labels. Each robustness row varies the label construction from the main \name configuration (marked with \colorbox{tahgray}{gray}); ``10\% noise'' randomly flips labels generated by the Qwen3-1.7B reference.}
\label{tab:oracle_label_robustness}
\centering
\footnotesize
\setlength{\tabcolsep}{6pt}
\begin{tabular}{llcccccc}
\toprule
\textbf{Method} & \textbf{Iter. Label} & \textbf{MATH500} & \textbf{GSM8K} & \textbf{AMC23} & \textbf{Olympiad} & \textbf{AIME25} & \textbf{Average} \\
\midrule
\multicolumn{8}{l}{\textit{Primary baselines}} \\
\quad Standard-1.7B & -- & 68.4 & 82.1 & 42.2 & 33.0 & 13.3 & 47.8 \\
\rowcolor{tahgray}
\quad \name-1.7B & 1.7B reference & \textbf{74.4} & \textbf{84.5} & \textbf{48.4} & \textbf{38.8} & \textbf{17.9} & \textbf{52.8}{\scriptsize$/+$5.0} \\
\addlinespace
\multicolumn{8}{l}{\textit{Label robustness checks}} \\
\quad \name-1.7B & 0.6B reference & 70.2 & 82.3 & 43.4 & 33.5 & 14.2 & 48.7{\scriptsize$/+$0.9} \\
\quad \name-1.7B & 1.7B reference + 10\% noise & \underline{73.6} & \underline{83.5} & \underline{46.8} & \underline{36.0} & \underline{14.6} & \underline{50.9}{\scriptsize$/+$3.1} \\
\bottomrule
\end{tabular}
\end{table*}

\begin{figure}[tb]
    \centering
    \includegraphics[width=0.8\columnwidth]{figure/flow.pdf}
    \caption{Next-token prediction changes across iterations. Top-2 frequent iterated tokens are visualized.}
    \label{fig:flow}
\end{figure}

% \xhdr{Generalizability}\todo{merge with performance; move to 5.3}
% We further study generalization when TaH is evaluated out of domain (OOD) or trained on broader data mixtures.
% First, when trained only on math data, TaH+ still improves OOD STEM performance on MMLU-STEM (4.7\% and 2.9\% for 0.6B and 1.7B respectively), indicating that the learned thinking patterns transfer robustly across domains.
% Second, fine-tuning Qwen3-1.7B-Base on a balanced OpenR1 mixture of math, QA, and code shows that TaH+ yields consistent gains over Standard across all categories, improving the overall average accuracy by 6.8\%.
% See additional experiment and performance details in Appendix~\ref{sec:appendix/exp/generalization}.

\xhdr{Continuation threshold}
\name shows robust performance across continuation thresholds and iteration ratios (Figure~\ref{fig:continuation_threshold}).
We empirically set $c_{\text{threshold}}=0.9$ for all evaluations.

\begin{table*}[t]
\caption{Performance comparison with larger maximum iteration depths. TaH variants are all 1.7B, trained on the math subset. Brackets show the percentage of tokens executing iterations 2, 3, and 4.}
\label{tab:max_interation_3}
\centering
\small
\setlength{\tabcolsep}{2pt}
\newcommand{\iterpct}[1]{\hspace{0.35em}{\scriptsize[#1]}}
\begin{tabular*}{\textwidth}{@{\extracolsep{\fill}}llcccccc@{}}
\toprule
\textbf{Method} & \textbf{Iter. Policy} & \textbf{MATH500} & \textbf{GSM8K} & \textbf{AMC23} & \textbf{Olympiad} & \textbf{AIME25} & \textbf{Avg.} \\
\midrule
Standard & -- & 68.4\iterpct{0.0, 0.0, 0.0} & 82.1\iterpct{0.0, 0.0, 0.0} & 42.2\iterpct{0.0, 0.0, 0.0} & 33.0\iterpct{0.0, 0.0, 0.0} & 13.3\iterpct{0.0, 0.0, 0.0} & 47.8\iterpct{0.0, 0.0, 0.0} \\
\midrule
TaH-2 & \multirow{3}{*}{Decreasing} & \underline{74.4}\iterpct{5.6, 0.0, 0.0} & \underline{84.5}\iterpct{7.5, 0.0, 0.0} & 48.4\iterpct{4.2, 0.0, 0.0} & 38.8\iterpct{5.7, 0.0, 0.0} & 17.9\iterpct{6.0, 0.0, 0.0} & 52.8\iterpct{5.8, 0.0, 0.0} \\
TaH-3 & & \textbf{74.8}\iterpct{6.5, 1.2, 0.0} & 84.0\iterpct{8.6, 1.8, 0.0} & \underline{49.1}\iterpct{6.3, 1.0, 0.0} & \textbf{41.6}\iterpct{5.3, 1.0, 0.0} & \underline{19.6}\iterpct{5.4, 1.0, 0.0} & \underline{53.8}\iterpct{6.4, 1.2, 0.0} \\
TaH-4 & & \textbf{74.8}\iterpct{4.9, 0.6, 0.6} & \textbf{84.6}\iterpct{7.1, 0.9, 0.8} & \textbf{49.7}\iterpct{4.3, 0.5, 0.4} & \underline{40.7}\iterpct{4.0, 0.5, 0.4} & \textbf{20.4}\iterpct{4.3, 0.5, 0.3} & \textbf{54.0}\iterpct{4.9, 0.6, 0.5} \\
\midrule
TaH-3 & Uniform & 70.2\iterpct{29.4, 31.8, 0.0} & 82.5\iterpct{29.5, 35.8, 0.0} & 41.6\iterpct{30.7, 32.9, 0.0} & 34.4\iterpct{30.6, 34.1, 0.0} & 14.2\iterpct{30.4, 35.3, 0.0} & 48.6\iterpct{30.1, 34.0, 0.0} \\
\bottomrule
\end{tabular*}
\end{table*}


\xhdr{Deeper iteration depth}
We train 1.7B variants with larger maximum iteration depths on the math subset.
TaH-3 and TaH-4 set $D_{\max}=3$ and $D_{\max}=4$, respectively, following the same sparse, decreasing allocation principle as TaH-2: most tokens stop at depth~1, while harder tokens continue through additional latent passes.
As shown in Table~\ref{tab:max_interation_3}, the average gain over Standard increases from +5.0\% (TaH-2) to +6.0\% (TaH-3) and +6.2\% (TaH-4) with little extra cost.
We further study the impact of the label distribution by comparing with TaH-3-Uniform, which assigns same number of tokens across depths.
The weaker result of the uniform variant confirms that sparser allocation of deeper iterations is more effective than forcing frequent use of deeper depths.
Detailed setup is provided in Appendix~\ref{sec:appendix/oracle_policy}.


\xhdr{Restricted inference depth}
We evaluate whether the benefits of \name persist when the maximum inference depth is below the training maximum.
Table~\ref{tab:restricted_inference_depth} evaluates 1.7B models trained with $D_{\max}=2$ but forced to verbalize after the first iteration at inference.
Although weaker than full \name inference, restricted-depth inference still outperforms Standard by 2.1-3.1\% on average, suggesting that \name training also improves first-iteration representations.

\begin{table}[tb]
    \centering
    \caption{Performance of 1.7B models under restricted inference depth, grouped by task category.}
    \label{tab:restricted_inference_depth}
    \small
    \setlength{\tabcolsep}{3pt}
    \begin{tabular}{lllcccc}
    \toprule
    \textbf{Method} & \makecell{\textbf{Train.}\\\textbf{Depth}} & \makecell{\textbf{Infer.}\\\textbf{Depth}} & \textbf{Math} & \textbf{Science} & \textbf{Code} & \textbf{Avg.} \\
    \midrule
    Standard & 1 & 1 & 46.5 & 52.2 & 45.5 & 47.5 \\
    \midrule
    \multirow{2}{*}{\name} & 2 & 2 & 49.6 & \underline{54.1} & \underline{52.7} & \underline{51.3}{\scriptsize$/+$3.8} \\
     & 2 & 1 & 48.1 & 53.9 & 49.2 & 49.6{\scriptsize$/+$2.1} \\
    \midrule
    \multirow{2}{*}{TaH+} & 2 & 2 & \textbf{51.7} & \textbf{58.0} & \textbf{54.5} & \textbf{53.7}{\scriptsize$/+$6.2} \\
     & 2 & 1 & \underline{50.2} & 52.4 & 49.8 & 50.6{\scriptsize$/+$3.1} \\
    \bottomrule
    \end{tabular}
\end{table}

%%%%%%%%%%%%%% Behavior %%%%%%%%%%%%%%%%%


% \begin{figure}[tb]
%     \centering
%     \begin{minipage}[b]{0.49\textwidth}
%         \centering
%         \begin{tabular}{l l c}
%         \toprule
%         \textbf{Training} & \textbf{Inference} & \textbf{Accuracy} \\
%         \midrule
%         Standard    & Standard    & 52 \\
%         \midrule
%         AlwaysThink & AlwaysThink & 38{\scriptsize$/-$14} \\
%         AlwaysThink & TaH-Oracle       & 77{\scriptsize$/+$25} \\
%         \midrule
%         TaH-Oracle   & TaH-Decider   & 54{\scriptsize$/+$\phantom{0}2} \\
%         TaH-Oracle   & TaH-Oracle       & 80{\scriptsize$/+$28} \\
%         \bottomrule
%         \end{tabular}
%         \captionof{table}{Impact of iteration strategies on Qwen3-0.6B (first 100 MATH500 samples).}
%         \label{tab:oracle}
%     \end{minipage}
%     \hfill
%     \begin{minipage}[b]{0.49\textwidth}
%         \centering
%         \includegraphics[width=\textwidth]{figure/pie.png}
%         \captionof{figure}{Next-token prediction accuracies on Qwen3-1.7B fine-tuned variants. Switching from direct reply to uniform think-twice causes more wrong corrections than fixes. Iterating only at hard tokens, identified by errors under the reference direct model, yields substantially more fixes with minimal harm.}
%         \label{fig:pie}
%     \end{minipage}
% \end{figure}
    
\subsection{Behavior Analysis}

\xhdr{Token alternation patterns}
We analyze which tokens \name selects for deeper iteration.
As shown in Figure~\ref{fig:flow}, \textit{But} and \textit{So} are iterated most frequently.
These tokens are hard to predict because they mark points of contrast or causality that can redirect subsequent reasoning.
At such junctures, \name uses additional iterations to refine its reasoning direction.
See Appendix~\ref{sec:appendix/exp/flow} for details.

\xhdr{Attention pattern}
We visualize the attention pattern of \name in Figure~\ref{fig:attention} and Appendix~\ref{sec:appendix/exp/attention}. Duo-causal attention focuses on different iterations in different heads, extracting broader contexts from multiple depths.

\begin{figure}[tb]
    \centering
    \includegraphics[width=0.67\columnwidth]{figure/landscape.png}
    \caption{Distribution of token prediction accuracy across iterations for Ouro and \name.}
    \label{fig:landscape}
\end{figure}

\xhdr{Accuracy landscape}
Figure~\ref{fig:landscape} shows NTP accuracy across iterations. Because Ouro trains all iterations to predict all tokens, predictable tokens across depths largely overlap. 
\name specializes deeper iterations for hard tokens, improving overall coverage under selective iteration.
